\chapter{Propósito}

\eyebrow{por que este software existe}

\vspace{4pt}
O Apex é uma aplicação de desktop para a \textbf{coleta e análise de imagens de
carcaças ovinas}, construída para montar um conjunto de dados de pesquisa e dar
continuidade à pesquisa de tipificação de carcaças no iCEV.

Seu propósito é restrito e deliberado. Um esforço anterior de coleta de dados não
fracassou por causa de seus métodos de aprendizado de máquina --- ele fracassou por
causa da \textbf{integridade do pareamento} entre cada imagem, o animal do qual ela
veio e o grau atribuído a ela. Esse vínculo não foi registrado no momento da coleta;
ele foi reconstruído posteriormente, pela ordem dos arquivos, sem verificação. As
etiquetas físicas estavam duplicadas ou ilegíveis, a planilha de laboratório não
tinha uma coluna com o nome da imagem, e apenas algumas etiquetas coincidiam com os
registros do laboratório. O resultado foi um conjunto de dados no qual nenhuma
imagem podia ser confiada como pertencente ao animal ao qual era atribuída.

\begin{note}[title={A ideia central}]
O Apex não é apenas um capturador de imagens. É a ferramenta que torna o pareamento
imagem\,$\leftrightarrow$\,animal\,$\leftrightarrow$\,grau \textbf{verificável
por construção}, de modo que o fracasso anterior não possa se repetir. Cada decisão
arquitetural serve a esse objetivo.
\end{note}

\begin{note}[title={Uma nota sobre o nome}]
A aplicação se chama \textbf{Apex}. Sua janela ainda carrega a marca interna
\tele{CARCASS} na barra de título e na barra de status (visível nas capturas de
tela), que é a assinatura do instrumento na tela; ao longo deste manual o produto é
referido como Apex.
\end{note}

\section{O que ele faz}

Em torno desse núcleo de integridade, o software fornece um fluxo de trabalho de
ponta a ponta:

\begin{itemize}[leftmargin=1.3em,itemsep=3pt]
  \item \textbf{Aquisição.} Registre uma carcaça com uma etiqueta física
        obrigatória, depois capture imagens a partir de uma câmera externa --- a
        imagem nasce pareada à carcaça no banco de dados, nunca inferida a partir de
        um nome de arquivo.
  \item \textbf{Importação.} Traga fotos capturadas em outro lugar; cada uma se
        torna uma carcaça (ou é reconciliada a uma existente), de modo que imagens
        externas entram no conjunto de dados sem quebrar o pareamento.
  \item \textbf{Análise.} Execute modelos treinados sobre as imagens: segmentação
        de gordura (validada), um grau de acabamento experimental e uma medida
        analítica de conformação (convexidade) --- tudo com remoção automática de
        fundo, para que as medições de superfície fiquem corretas.
  \item \textbf{Avaliação.} Colete graus independentes e cegos de múltiplos
        avaliadores e meça a concordância entre eles ($\kappa$ de Fleiss).
  \item \textbf{Monitor ao vivo.} Uma sobreposição de gordura em tempo real a partir
        de uma câmera ou de um arquivo de vídeo, para ajudar a enquadrar a carcaça
        durante a captura.
  \item \textbf{Painel \& exportação.} Acompanhe o progresso da amostra em relação
        à meta e exporte o conjunto de dados como um manifesto mais um relatório de
        integridade.
\end{itemize}

\section{Para quem ele é}

Os usuários principais são a equipe de pesquisa e os operadores que coletam imagens
de carcaças na linha de abate e no laboratório. A interface é projetada como um
\emph{instrumento} científico --- examinado e operado, não lido de cima a baixo ---
e segue a linguagem visual do software de controle de missão: um fundo escuro de
instrumento, um destaque de sinal ciano e leituras em fonte monoespaçada para cada
valor medido.

\section{Uma nota sobre honestidade científica}

Dois dos números que este software produz --- o \emph{grau de acabamento} e o
\emph{grau de conformação} --- estão marcados em toda a interface como
\textbf{experimentais / não validados}. Esta é uma escolha deliberada, fundamentada
nos próprios achados do projeto (Estação~13). O software os exibe porque são úteis
como referências exploratórias e porque o conjunto de dados que ele constrói é
exatamente o que se precisa para validá-los no futuro. Ele nunca os apresenta como
fatos consolidados.
